% preambulo.tex — Estética de los informes de laboratorio de Física (PDF).
% Replica la de los informes TP5/TP6 (A1193-TP5-G5.pdf): Liberation Sans
% (métrica de Arial), títulos azul marino, subtítulos azules, encabezado con
% filete azul, tablas de encabezado azul. Lo incluye lab_fisica/pdf.yaml.

\usepackage{tikz}
\usetikzlibrary{babel}   % babel spanish activa < > como atajos: sin esto, TikZ falla
\usepackage{float}
\floatplacement{figure}{H}
\usepackage{graphicx}

% Colores tomados del TP5 (los mismos que usa docx_lab.py)
\definecolor{azullab}{HTML}{2E75B6}
\definecolor{marinolab}{HTML}{1F3864}
\definecolor{grislab}{HTML}{676767}
\definecolor{bordelab}{HTML}{AAAAAA}
\definecolor{grisfondo}{HTML}{E6E6E6}

% Carátula: Times (Liberation Serif, misma métrica que Times New Roman)
\newfontfamily\caratulafuente{Liberation Serif}[Ligatures=TeX]
\newcommand{\caratulagrande}{\caratulafuente\fontsize{16}{19}\selectfont\bfseries}
\newcommand{\caratulanombres}{\caratulafuente\fontsize{16}{19}\selectfont\mdseries}
\newcommand{\caratulachica}{\caratulafuente\fontsize{11}{13}\selectfont\bfseries}
% Texto de una línea que se achica si no entra en el ancho dado
\newsavebox{\labcaja}
\newcommand{\labajustar}[2]{\sbox{\labcaja}{#2}\ifdim\wd\labcaja>#1\resizebox{#1}{!}{\usebox{\labcaja}}\else\usebox{\labcaja}\fi}

% Cuerpo
\linespread{1.1}
\usepackage{titlesec}
\titleformat{\section}{\Large\bfseries\color{marinolab}}{}{0pt}{}
\titleformat{\subsection}{\large\bfseries\color{azullab}}{}{0pt}{}
\titleformat{\subsubsection}{\normalsize\bfseries\color{azullab}}{}{0pt}{}
\titlespacing*{\section}{0pt}{16pt}{6pt}
\titlespacing*{\subsection}{0pt}{12pt}{4pt}
\titlespacing*{\subsubsection}{0pt}{10pt}{2pt}

% Encabezado y pie: el texto del encabezado lo fija lab-fisica.lua
\usepackage{fancyhdr}
\pagestyle{fancy}
\fancyhf{}
\fancyfoot[R]{\small\thepage}
\renewcommand{\headrule}{\vspace{-2pt}{\color{azullab}\hrule height 0.6pt}}
\renewcommand{\footrule}{{\color{azullab}\hrule height 0.6pt}\vspace{4pt}}
\setlength{\headheight}{14pt}

% Pies de figura (los de tabla los arma lab-fisica.lua)
\usepackage{caption}
\captionsetup{font={footnotesize,color=grislab},labelsep=colon}

% Tablas: texto chico, filas altas, grilla gris; poco aire antes del pie
\newcommand{\tablalab}{\small\renewcommand{\arraystretch}{1.5}\arrayrulecolor{bordelab}}
\setlength{\LTpost}{2pt}

% Recuadro de resultado: filete azul, a todo el ancho, contenido centrado
\usepackage[most]{tcolorbox}
\newtcolorbox{resultadolab}{sharp corners, colback=white, colframe=azullab,
  boxrule=0.6pt, left=6pt, right=6pt, top=4pt, bottom=4pt,
  before skip=8pt, after skip=8pt, halign=flush center}

% Caracteres Unicode sin glifo en Liberation Sans
\usepackage{newunicodechar}
\newunicodechar{≈}{\ensuremath{\approx}}
